% generated by bibm/supplement/make_sm_tables.py -- do not edit; sources: gue_twins.csv
\begin{table}[!htbp]
\caption{GUE mouse-TF split artifact: dinucleotide-shuffle negatives were split independently of their source positives (post hoc PH2). Source: \src{gue\_twins.csv}.}
\label{tab:guetwins}
\centering
\footnotesize
\begin{tabular}{lrrrrr}
\toprule
Dataset & Test $n$ & Test w/ cross-label twin in train & Comp.-only AUROC, all & with twin & without twin \\
\midrule
mouse\_tf\_0 & 810 & 81.9\% & 0.023 & 0.000 & 0.520 \\
mouse\_tf\_1 & 6{,}745 & 80.5\% & 0.081 & 0.015 & 0.516 \\
mouse\_tf\_2 & 328 & 81.1\% & 0.018 & 0.000 & 0.493 \\
mouse\_tf\_3 & 239 & 72.8\% & 0.033 & 0.000 & 0.452 \\
mouse\_tf\_4 & 1{,}883 & 80.6\% & 0.034 & 0.000 & 0.458 \\
\bottomrule
\end{tabular}
\par\smallskip
\begin{minipage}{\linewidth}\footnotesize\raggedright
A test sequence whose shuffle twin (same composition, opposite label) is in the training split inverts a composition-only model; on test sequences without such a twin the composition-only AUROC is close to 0.5. When negatives are shuffles of the positives, each positive and its shuffle belong in the same split.
\end{minipage}
\end{table}
